%------------------------------------------------------------------------------%
\documentclass[fleqn,utf8,aspectratio=169,12pt]{beamer}

\usepackage{gaal}

\title{Produto Misto}

%------------------------------------------------------------------------------%
\begin{document}

\addtitle

%------------------------------------------------------------------------------%
\section{Definição}

%------------------------------------------------------------------------------%
\begin{frame}{Produto Misto}
  Dados
  \[
    \vec{v}, \, \vec{w}, \, \vec{z}\in\R^3
  \]

  \pause
  Produto misto
  \[
    \left(\vec{v}\times\vec{w}\right)\cdot\vec{z}
  \]

  \pause
  O produto misto é um escalar
  \[
  \left(\vec{v}\times\vec{w}\right)\cdot\vec{z} \in \R
  \]
\end{frame}

%------------------------------------------------------------------------------%
%------------------------------------------------------------------------------%
\begin{question}
  Calcule o produto misto dos vetores
  \[
    \vec{v} = \vetortri{1}{ 2}{ 3}
    \qquad
    \vec{w} = \vetortri{2}{-1}{ 1}
    \qquad
    \vec{z} = \vetortri{3}{ 1}{-2}
  \]
\end{question}

%------------------------------------------------------------------------------%
\begin{resolution}{Resolução}
  \begin{align*}
    \onslide<+->{\vec{v}\times\vec{w}}
    \onslide<+->{& = \begin{vmatrix}
      \vec{i} & \vec{j} & \vec{k} \\
      1       & 2       & 3       \\
      2       & -1      & 1
    \end{vmatrix}}
    \\
    \onslide<+->{& = \vec{i} (-1)^{1+1}
        \begin{vmatrix}
          2 & 3 \\
          -1 & 1
        \end{vmatrix}
      \;+\; \vec{j} (-1)^{1+2}
        \begin{vmatrix}
          1 & 3 \\
          2 & 1
        \end{vmatrix}
      \;+\; \vec{k} (-1)^{1+3}
        \begin{vmatrix}
          1 & 2 \\
          2 & -1
        \end{vmatrix}}
    \\
    \onslide<+->{& = \vec{i} (2+3)
      - \vec{j} (1-6)
      + \vec{k} (-1-4)}
    \\
    \onslide<+->{& = 5 \vec{i} + 5 \vec{j} - 5 \vec{k}}
    \\
    \onslide<+->{& = \vetortri{5}{5}{-5}}
  \end{align*}
\end{resolution}

%------------------------------------------------------------------------------%
\begin{resolution}{Resolução}
  \[
    \left(\vec{v}\times\vec{w}\right)\cdot\vec{z}
    \pause
    = \vetortri{5}{5}{-5} \cdot \vetortri{3}{ 1}{-2}
    \pause
    = 15 + 5 + 10
    \pause
    = 30
  \]
\end{resolution}

%------------------------------------------------------------------------------%


%------------------------------------------------------------------------------%
\begin{frame}{Calculando o Produto Misto Diretamente}
  Se
  \[
    \vec{v} = \vetortri{v_1}{v_2}{v_3}
    \qquad
    \vec{w} = \vetortri{w_1}{w_2}{w_3}
    \qquad
    \vec{z} = \vetortri{z_1}{z_2}{z_3}
  \]
  \pause
  então
  \[
    \left(\vec{v}\times\vec{w}\right)\cdot\vec{z}
    \pause
    = \begin{vmatrix}
      v_1 & v_2 & v_3 \\
      w_1 & w_2 & w_3 \\
      z_1 & z_2 & z_3
    \end{vmatrix}
  \]
\end{frame}

%------------------------------------------------------------------------------%
%------------------------------------------------------------------------------%
\begin{question}
Calcule o produto misto
\(
  \left(\vec{v}\times\vec{w}\right)\cdot\vec{z}
\)
para
\[
  \vec{v} = \vetortri{2}{ 1}{0}
  \qquad
  \vec{w} = \vetortri{1}{-1}{2}
  \qquad
  \vec{z} = \vetortri{3}{ 2}{1}
\]
\end{question}

%------------------------------------------------------------------------------%
\begin{resolution}{Solução}
  \begin{align*}
  \onslide<+->{\left(\vec{v}\times\vec{w}\right)\cdot\vec{z}}
  \onslide<+->{& = \begin{vmatrix}
      2 & 1  & 0 \\
      1 & -1 & 2 \\
      3 & 2  & 1
    \end{vmatrix}}
  \\
  \onslide<+->{& = 2
    \begin{vmatrix}
      -1 & 2 \\
      2  & 1
    \end{vmatrix}
    - \begin{vmatrix}
      1 & 2 \\
      3 & 1
    \end{vmatrix}
    + 0}
  \\
  \onslide<+->{& = 2(-1-4)-(1-6)}
  \\
  \onslide<+->{& = 2(-5)-(-5)}
  \\
  \onslide<+->{& = - 5}
  \end{align*}
\end{resolution}



%------------------------------------------------------------------------------%
\section{Propriedades}

%------------------------------------------------------------------------------%
\begin{frame}{Sinal do Produto Misto}
  Para
  \[
    p = \left(\vec{v}\times\vec{w}\right)\cdot\vec{z}
  \]

  \pause
  Se $\vec{z}$ aponta para o lado de $\vec{v}\times\vec{w}$
  \[
    p > 0
  \]

  \pause
  Se $\vec{z}$ aponta para o lado oposta a $\vec{v}\times\vec{w}$
  \[
    p < 0
  \]

  \pause
  Se $\vec{z}$ é paralelo ao plano determinado por $\vec{v}$ e $\vec{w}$
  \[
    p = 0
  \]
\end{frame}

%------------------------------------------------------------------------------%
\begin{frame}{Propriedade de Permutação}
  \onslide<+->{\emph{Trocas} de duas posições mudam o sinal}
  \begin{align*}
    \onslide<+->{\left(\vec{v}\times\vec{w}\right)\cdot\vec{z} & = - \left(\vec{w}\times\vec{v}\right)\cdot\vec{z} \\}
    \onslide<+->{\left(\vec{v}\times\vec{w}\right)\cdot\vec{z} & = - \left(\vec{v}\times\vec{z}\right)\cdot\vec{w} \\}
    \onslide<+->{\left(\vec{v}\times\vec{w}\right)\cdot\vec{z} & = - \left(\vec{z}\times\vec{w}\right)\cdot\vec{v}}
  \end{align*}

  \onslide<+->{\emph{Permutações cíclicas} preservam o valor}
  \[
      \onslide<+->{\left(\vec{v}\times\vec{w}\right)\cdot\vec{z}}
      \onslide<+->{\;=\; \left(\vec{w}\times\vec{z}\right)\cdot\vec{v}}
      \onslide<+->{\;=\; \left(\vec{z}\times\vec{v}\right)\cdot\vec{w}}
  \]
\end{frame}

%------------------------------------------------------------------------------%
%------------------------------------------------------------------------------%
\begin{question}
  Dados
  \[
    \vec{v} = \vetortri{1}{ 2}{0}
    \qquad
    \vec{w} = \vetortri{2}{-1}{1}
    \qquad
    \vec{z} = \vetortri{1}{ 0}{3}
  \]
  Calcule
  \(
    \left(\vec{v}\times\vec{w}\right)\cdot\vec{z}
  \) e
  \(
    \left(\vec{w}\times\vec{v}\right)\cdot\vec{z}
  \)
\end{question}

%------------------------------------------------------------------------------%
\begin{resolution}{Resolução}
  \begin{align*}
    \onslide<+->{\left(\vec{v}\times\vec{w}\right)\cdot\vec{z}}
    \onslide<+->{& = \begin{vmatrix}
      1 & 2  & 0 \\
      2 & -1 & 1 \\
      1 & 0  & 3
    \end{vmatrix}}
    \\
    \onslide<+->{& = 1 \begin{vmatrix}
        -1 & 1 \\
        0  & 3
      \end{vmatrix}
      - 2
      \begin{vmatrix}
        2 & 1 \\
        1 & 3
      \end{vmatrix}}
    \\
    \onslide<+->{& = -3 - 10}
    \\
    \onslide<+->{& = -13}
  \end{align*}

  \[
    \onslide<+->{\left(\vec{w}\times\vec{v}\right)\cdot\vec{z}}
    \onslide<+->{= -\left(\vec{v}\times\vec{w}\right)\cdot\vec{z}}
    \onslide<+->{= - (-13)}
    \onslide<+->{= 13}
  \]
\end{resolution}

%------------------------------------------------------------------------------%


%------------------------------------------------------------------------------%
\section{Coplanaridade}

%------------------------------------------------------------------------------%
\begin{frame}{Coplanaridade}
  Três vetores de $\R^3$ são \emph{coplanares} quando pertencem a um mesmo plano

  \pause
  \bigskip
  Condição equivalente
  \[
    \left(\vec{v}\times\vec{w}\right)\cdot\vec{z} = 0
  \]
\end{frame}

%------------------------------------------------------------------------------%
%------------------------------------------------------------------------------%
\begin{question}
  Verifique se os vetores são coplanares
  \[
    \vec{v} = \vetortri{1}{ 2}{1}
    \qquad
    \vec{w} = \vetortri{2}{-1}{3}
    \qquad
    \vec{z} = \vetortri{5}{ 3}{5}
  \]

  \pause
  \bigskip
  Condição de coplanaridade
  \[
    \left(\vec{v}\times\vec{w}\right)\cdot\vec{z} = 0
  \]
\end{question}

%------------------------------------------------------------------------------%
\begin{resolution}{Resolução}
  \begin{align*}
    \onslide<+->{\left(\vec{v}\times\vec{w}\right)\cdot\vec{z}}
    \onslide<+->{& = \begin{vmatrix}
        1 & 2  & 1 \\
        2 & -1 & 3 \\
        5 & 3  & 5
      \end{vmatrix}}
    \\
    \onslide<+->{& = \begin{vmatrix}
        -1 & 3 \\
        3  & 5
      \end{vmatrix}
      - 2 \begin{vmatrix}
        2 & 3 \\
        5 & 5
      \end{vmatrix}
      + \begin{vmatrix}
          2 & -1 \\
          5 & 3
        \end{vmatrix}}
    \\
    \onslide<+->{& = (-14) -2(-5) + 11 \\}
    \onslide<+->{& = 7}
    \onslide<+->{\neq 0}
  \end{align*}
  \onslide<+->{Os três vetores não são coplanares}
\end{resolution}

%------------------------------------------------------------------------------%
\begin{question}
  Determine $a$ para que os vetores sejam coplanares
  \[
    \vec{v} = \vetortri{1}{ 2}{3}
    \qquad
    \vec{w} = \vetortri{2}{-1}{1}
    \qquad
    \vec{z} = \vetortri{a}{ 3}{4}
  \]

  \pause
  \bigskip
  Condição de coplanaridade
  \[
    \left(\vec{v}\times\vec{w}\right)\cdot\vec{z} = 0
  \]
\end{question}

%------------------------------------------------------------------------------%
\begin{resolution}{Resolução}
\vspace{-1.5\baselineskip}
\begin{columns}[t]
\column{0.6\textwidth}
  \begin{align*}
  \onslide<+->{p & = \left(\vec{v}\times\vec{w}\right)\cdot\vec{z} \\}
  \onslide<+->{& = \begin{vmatrix}
    1 & 2  & 3 \\
    2 & -1 & 1 \\
    a & 3  & 4
  \end{vmatrix}}
  \\
  \onslide<+->{& = \begin{vmatrix}
      -1 & 1 \\
      3  & 4
    \end{vmatrix}
    - 2 \begin{vmatrix}
      2 & 1 \\
      a & 4
    \end{vmatrix}
    + 3 \begin{vmatrix}
      2 & -1 \\
      a & 3
    \end{vmatrix}}
  \\
  \onslide<+->{& = -7 - 2(8-a) + 3(6+a)   \\}
  \onslide<+->{& = -7 - 16 + 2a + 18 + 3a \\}
  \onslide<+->{& = 5a - 5}
  \end{align*}
\column{0.4\textwidth}
  \begin{align*}
    \onslide<+->{p    & = 0 \\}
    \onslide<+->{5a-5 & = 0 \\}
    \onslide<+->{5a   & = 5 \\}
    \onslide<+->{a    & = 1}
  \end{align*}
\end{columns}
\end{resolution}



%%------------------------------------------------------------------------------%
%\section{Lista Mínima}
%
%%------------------------------------------------------------------------------%
%\begin{frame}{Lista Mínima}
%  \begin{enumerate}
%    \item
%  \end{enumerate}
%
%  \vfill
%  Atenção: A prova é baseada no livro, não nas apresentações
%\end{frame}

%------------------------------------------------------------------------------%
\end{document}
%------------------------------------------------------------------------------%
